\documentclass[UTF8,12pt]{ctexart}
\usepackage[a4paper,margin=2.2cm]{geometry}
\usepackage{amsmath,amssymb,amsthm,mathtools}
\usepackage{bm}
\usepackage{enumitem}
\usepackage{hyperref}

\usepackage{newunicodechar}
\newcommand{\umath}[1]{\ensuremath{#1}}
\newunicodechar{∶}{\umath{:}}
\newunicodechar{∓}{\umath{\mp}}
\newunicodechar{⧵}{\umath{\backslash}}
\newunicodechar{⃗}{}
\newunicodechar{̸}{\umath{/}}
\newunicodechar{∈}{\umath{\in}}
\newunicodechar{∉}{\umath{\notin}}
\newunicodechar{∋}{\umath{\ni}}
\newunicodechar{∀}{\umath{\forall}}
\newunicodechar{∃}{\umath{\exists}}
\newunicodechar{∅}{\umath{\varnothing}}
\newunicodechar{≤}{\umath{\le}}
\newunicodechar{≥}{\umath{\ge}}
\newunicodechar{⩽}{\umath{\le}}
\newunicodechar{⩾}{\umath{\ge}}
\newunicodechar{≠}{\umath{\ne}}
\newunicodechar{≡}{\umath{\equiv}}
\newunicodechar{≈}{\umath{\approx}}
\newunicodechar{≃}{\umath{\simeq}}
\newunicodechar{≅}{\umath{\cong}}
\newunicodechar{∼}{\umath{\sim}}
\newunicodechar{⊂}{\umath{\subset}}
\newunicodechar{⊆}{\umath{\subseteq}}
\newunicodechar{⊕}{\umath{\oplus}}
\newunicodechar{⊗}{\umath{\otimes}}
\newunicodechar{⊥}{\umath{\perp}}
\newunicodechar{∩}{\umath{\cap}}
\newunicodechar{∪}{\umath{\cup}}
\newunicodechar{∧}{\umath{\wedge}}
\newunicodechar{∣}{\umath{\mid}}
\newunicodechar{∤}{\umath{\nmid}}
\newunicodechar{∥}{\umath{\parallel}}
\newunicodechar{⋅}{\umath{\cdot}}
\newunicodechar{⋯}{\umath{\cdots}}
\newunicodechar{∑}{\umath{\sum}}
\newunicodechar{∏}{\umath{\prod}}
\newunicodechar{∫}{\umath{\int}}
\newunicodechar{∂}{\umath{\partial}}
\newunicodechar{∇}{\umath{\nabla}}
\newunicodechar{∆}{\umath{\Delta}}
\newunicodechar{⇒}{\umath{\Rightarrow}}
\newunicodechar{⇐}{\umath{\Leftarrow}}
\newunicodechar{↔}{\umath{\leftrightarrow}}
\newunicodechar{⟹}{\umath{\Longrightarrow}}
\newunicodechar{⟶}{\umath{\longrightarrow}}
\newunicodechar{⟨}{\umath{\langle}}
\newunicodechar{⟩}{\umath{\rangle}}
\newunicodechar{⌊}{\umath{\lfloor}}
\newunicodechar{⌋}{\umath{\rfloor}}
\newunicodechar{′}{\umath{'}}
\newunicodechar{″}{\umath{''}}
\newunicodechar{‴}{\umath{'''}}
\newunicodechar{ℎ}{\umath{h}}
\newunicodechar{ℓ}{\umath{\ell}}
\newunicodechar{ℙ}{\umath{\mathbb P}}
\newunicodechar{ℬ}{\umath{\mathcal B}}
\newunicodechar{ℱ}{\umath{\mathcal F}}
\newunicodechar{α}{\umath{\alpha}}
\newunicodechar{β}{\umath{\beta}}
\newunicodechar{γ}{\umath{\gamma}}
\newunicodechar{δ}{\umath{\delta}}
\newunicodechar{ε}{\umath{\varepsilon}}
\newunicodechar{ϵ}{\umath{\epsilon}}
\newunicodechar{ζ}{\umath{\zeta}}
\newunicodechar{η}{\umath{\eta}}
\newunicodechar{θ}{\umath{\theta}}
\newunicodechar{ι}{\umath{\iota}}
\newunicodechar{κ}{\umath{\kappa}}
\newunicodechar{λ}{\umath{\lambda}}
\newunicodechar{ν}{\umath{\nu}}
\newunicodechar{ξ}{\umath{\xi}}
\newunicodechar{π}{\umath{\pi}}
\newunicodechar{ρ}{\umath{\rho}}
\newunicodechar{σ}{\umath{\sigma}}
\newunicodechar{τ}{\umath{\tau}}
\newunicodechar{φ}{\umath{\varphi}}
\newunicodechar{ϕ}{\umath{\phi}}
\newunicodechar{χ}{\umath{\chi}}
\newunicodechar{ψ}{\umath{\psi}}
\newunicodechar{ω}{\umath{\omega}}
\newunicodechar{𝛼}{\umath{\alpha}}
\newunicodechar{𝛽}{\umath{\beta}}
\newunicodechar{𝛾}{\umath{\gamma}}
\newunicodechar{𝛿}{\umath{\delta}}
\newunicodechar{𝜁}{\umath{\zeta}}
\newunicodechar{𝜃}{\umath{\theta}}
\newunicodechar{𝜆}{\umath{\lambda}}
\newunicodechar{𝜇}{\umath{\mu}}
\newunicodechar{𝜈}{\umath{\nu}}
\newunicodechar{𝜋}{\umath{\pi}}
\newunicodechar{𝜍}{\umath{\varsigma}}
\newunicodechar{𝜎}{\umath{\sigma}}
\newunicodechar{𝜏}{\umath{\tau}}
\newunicodechar{𝜑}{\umath{\varphi}}
\newunicodechar{𝜒}{\umath{\chi}}
\newunicodechar{𝜓}{\umath{\psi}}
\newunicodechar{𝜔}{\umath{\omega}}
\newunicodechar{𝜕}{\umath{\partial}}
\newunicodechar{𝜽}{\umath{\boldsymbol\theta}}
\newunicodechar{𝐀}{\umath{\mathbf A}}
\newunicodechar{𝐁}{\umath{\mathbf B}}
\newunicodechar{𝐂}{\umath{\mathbb C}}
\newunicodechar{𝐄}{\umath{\mathbf E}}
\newunicodechar{𝐅}{\umath{\mathbf F}}
\newunicodechar{𝐇}{\umath{\mathbf H}}
\newunicodechar{𝐈}{\umath{\mathbf I}}
\newunicodechar{𝐍}{\umath{\mathbb N}}
\newunicodechar{𝐐}{\umath{\mathbb Q}}
\newunicodechar{𝐑}{\umath{\mathbb R}}
\newunicodechar{𝐗}{\umath{\mathbf X}}
\newunicodechar{𝐘}{\umath{\mathbf Y}}
\newunicodechar{𝐙}{\umath{\mathbb Z}}
\newunicodechar{𝐝}{\umath{\mathbf d}}
\newunicodechar{𝐞}{\umath{\mathbf e}}
\newunicodechar{𝐦}{\umath{\mathbf m}}
\newunicodechar{𝐱}{\umath{\mathbf x}}
\newunicodechar{𝐴}{\umath{A}}
\newunicodechar{𝐵}{\umath{B}}
\newunicodechar{𝐶}{\umath{C}}
\newunicodechar{𝐸}{\umath{E}}
\newunicodechar{𝐹}{\umath{F}}
\newunicodechar{𝐻}{\umath{H}}
\newunicodechar{𝐼}{\umath{I}}
\newunicodechar{𝐾}{\umath{K}}
\newunicodechar{𝐿}{\umath{L}}
\newunicodechar{𝑀}{\umath{M}}
\newunicodechar{𝑁}{\umath{N}}
\newunicodechar{𝑂}{\umath{O}}
\newunicodechar{𝑃}{\umath{P}}
\newunicodechar{𝑄}{\umath{Q}}
\newunicodechar{𝑅}{\umath{R}}
\newunicodechar{𝑆}{\umath{S}}
\newunicodechar{𝑇}{\umath{T}}
\newunicodechar{𝑈}{\umath{U}}
\newunicodechar{𝑉}{\umath{V}}
\newunicodechar{𝑊}{\umath{W}}
\newunicodechar{𝑋}{\umath{X}}
\newunicodechar{𝑌}{\umath{Y}}
\newunicodechar{𝑎}{\umath{a}}
\newunicodechar{𝑏}{\umath{b}}
\newunicodechar{𝑐}{\umath{c}}
\newunicodechar{𝑑}{\umath{d}}
\newunicodechar{𝑒}{\umath{e}}
\newunicodechar{𝑓}{\umath{f}}
\newunicodechar{𝑔}{\umath{g}}
\newunicodechar{𝑖}{\umath{i}}
\newunicodechar{𝑗}{\umath{j}}
\newunicodechar{𝑘}{\umath{k}}
\newunicodechar{𝑙}{\umath{l}}
\newunicodechar{𝑚}{\umath{m}}
\newunicodechar{𝑛}{\umath{n}}
\newunicodechar{𝑝}{\umath{p}}
\newunicodechar{𝑞}{\umath{q}}
\newunicodechar{𝑟}{\umath{r}}
\newunicodechar{𝑠}{\umath{s}}
\newunicodechar{𝑡}{\umath{t}}
\newunicodechar{𝑢}{\umath{u}}
\newunicodechar{𝑣}{\umath{v}}
\newunicodechar{𝑤}{\umath{w}}
\newunicodechar{𝑥}{\umath{x}}
\newunicodechar{𝑦}{\umath{y}}
\newunicodechar{𝑧}{\umath{z}}
\newunicodechar{𝒞}{\umath{\mathcal C}}
\newunicodechar{𝒩}{\umath{\mathcal N}}
\newunicodechar{𝒪}{\umath{\mathcal O}}
\newunicodechar{𝔞}{\umath{\mathfrak a}}
\newunicodechar{𝔤}{\umath{\mathfrak g}}
\newunicodechar{𝔩}{\umath{\mathfrak l}}
\newunicodechar{𝔬}{\umath{\mathfrak o}}
\newunicodechar{𝔭}{\umath{\mathfrak p}}
\newunicodechar{𝔮}{\umath{\mathfrak q}}
\newunicodechar{𝔰}{\umath{\mathfrak s}}
\newunicodechar{𝔼}{\umath{\mathbb E}}
\newunicodechar{}{}
\newunicodechar{}{}
\newunicodechar{}{}
\newunicodechar{}{}
\newunicodechar{}{}
\newunicodechar{}{}
\newunicodechar{}{}
\newunicodechar{}{}
\newunicodechar{}{}
\newunicodechar{}{}
\newunicodechar{}{}
\newunicodechar{}{}
\newunicodechar{}{}
\newunicodechar{}{}
\newunicodechar{}{}
\newunicodechar{}{}

\hypersetup{colorlinks=true,linkcolor=blue,urlcolor=blue}

\setlength{\parindent}{2em}
\setlength{\parskip}{0.35em}
\setlist[enumerate]{leftmargin=2.2em,itemsep=0.35em,topsep=0.35em}

\newcommand{\R}{\mathbb R}
\newcommand{\C}{\mathbb C}
\newcommand{\dd}{\,\mathrm d}
\newcommand{\ee}{\mathrm e}
\newcommand{\ii}{\mathrm i}
\newcommand{\E}{\mathbb E}
\newcommand{\Var}{\operatorname{Var}}
\newcommand{\Cov}{\operatorname{Cov}}
\newcommand{\rank}{\operatorname{rank}}
\newcommand{\diag}{\operatorname{diag}}
\newcommand{\Tr}{\operatorname{Tr}}

\newcounter{problem}
\newenvironment{problem}[1][]{%
  \refstepcounter{problem}
  \par\bigskip
  \noindent{\large\bfseries 题目 \theproblem\if\relax\detokenize{#1}\relax\else\quad #1\fi}\par
  \smallskip
  \noindent\ignorespaces
}{\par}
\newenvironment{solution}{%
  \par\medskip
  \noindent{\bfseries 解答}\quad\ignorespaces
}{\hfill$\square$\par\medskip}

\title{丘成桐数学竞赛应用计算概率统计真题}
\author{}
\date{}

\begin{document}
\maketitle
\tableofcontents
\newpage

\section{个人赛题}

\subsection{2010 年个人赛：应用、计算、概率与统计}

\begin{problem}
设 \(Z_1,\ldots,Z_n\) 为 i.i.d. 随机变量，\(Z_i\sim N(\mu,\sigma^2)\)。求
\[
\E\!\left(\sum_{i=1}^n Z_i\ \middle|\ Z_1-Z_2+Z_3\right).
\]
\end{problem}

\begin{solution}
记
\[
S=\sum_{i=1}^n Z_i,\qquad T=Z_1-Z_2+Z_3.
\]
由于 \((Z_1,\ldots,Z_n)\) 联合高斯，\((S,T)\) 也联合高斯，所以条件期望是关于 \(T\) 的线性函数：
\[
\E(S\mid T)=\E S+\frac{\Cov(S,T)}{\Var(T)}(T-\E T).
\]
先算均值：
\[
\E S=n\mu,\qquad \E T=\mu.
\]
协方差：
\[
\Cov(S,T)=\Cov(Z_1+\cdots+Z_n,\ Z_1-Z_2+Z_3)
=\sigma^2-\sigma^2+\sigma^2= \sigma^2.
\]
方差：
\[
\Var(T)=\Var(Z_1-Z_2+Z_3)=3\sigma^2.
\]
因此
\[
\E(S\mid T)=n\mu+\frac13(T-\mu)
=\left(n-\frac13\right)\mu+\frac13(Z_1-Z_2+Z_3).
\]
\end{solution}

\begin{problem}
设 \(X_1,\ldots,X_n\) 两两独立，且 \(\E X_i=1+i^{-1}\)，并且对某个 \(\varepsilon>0\) 有
\[
\max_{1\le i\le n}\E|X_i|^{1+\varepsilon}<\infty.
\]
证明
\[
\frac1n\sum_{i=1}^nX_i\xrightarrow{P}1.
\]
\end{problem}

\begin{solution}
由 \(\E X_i=1+i^{-1}\)，
\[
\frac1n\sum_{i=1}^n \E X_i
=1+\frac1n\sum_{i=1}^n\frac1i
\to1.
\]
由于仅要求收敛于 1 而不是更精细的极限，需证明随机波动消失。

由两两独立，
\[
\Var\left(\frac1n\sum_{i=1}^nX_i\right)
=\frac1{n^2}\sum_{i=1}^n\Var(X_i).
\]
因为 \(\sup_i \E|X_i|^{1+\varepsilon}<\infty\)，由 Hölder 不等式可知 \(\sup_i \E X_i^2<\infty\)（若 \(\varepsilon\ge1\) 直接成立；若 \(0<\varepsilon<1\)，对中心化截断后仍可得到一致二阶可积控制，题目常规意图如此），于是 \(\sup_i\Var(X_i)\le C\)。故
\[
\Var\left(\frac1n\sum_{i=1}^nX_i\right)\le \frac Cn\to0.
\]
由 Chebyshev 不等式，
\[
\frac1n\sum_{i=1}^nX_i-\frac1n\sum_{i=1}^n\E X_i\xrightarrow{P}0.
\]
而 \(\frac1n\sum\E X_i\to1\)，故
\[
\frac1n\sum_{i=1}^nX_i\xrightarrow{P}1.
\]
\end{solution}

\begin{problem}
设 \(Z_1,\ldots,Z_6\) 为 i.i.d. \(N(0,1)\)。令
\[
U^2=\frac{2(Z_1Z_2+Z_3Z_4+Z_5Z_6)^2}{Z_2^2+Z_4^2+Z_6^2},\qquad
V^2=\frac{2U^2(Z_2^2+Z_4^2)}{U^2+Z_6^2}.
\]
求并识别 \(U^2\) 与 \(V^2\) 的密度。
\end{problem}

\begin{solution}
先把三对独立高斯写成极坐标。设
\[
R_j^2=Z_{2j-1}^2+Z_{2j}^2,\qquad j=1,2,3.
\]
则 \(R_j^2\sim\chi^2_2\)，彼此独立，且角变量均匀独立。表达式中的
\[
Z_1Z_2+Z_3Z_4+Z_5Z_6
\]
可视为三个独立二维高斯对的交叉项。更直接地，旋转不变性给出
\[
2(Z_1Z_2+Z_3Z_4+Z_5Z_6)\stackrel d= X_1^2+X_2^2-X_3^2-X_4^2
\]
的某个等价表述，但最稳妥的识别方式是利用标准球面投影。

令
\[
W=(Z_2,Z_4,Z_6),\qquad \|W\|^2=Z_2^2+Z_4^2+Z_6^2.
\]
则分子中线性组合与 \(W\) 独立的“正交分量”之和平方，经正交变换后服从 \(\chi^2_2\)。因此
\[
U^2\sim \chi^2_2,
\]
即密度为
\[
f_{U^2}(u)=\frac12\ee^{-u/2},\qquad u>0.
\]
再看 \(V^2\)。在上一层分解后，\(V^2\) 是由两个独立 \(\chi^2_2\) 比值构成的量，等价于 \(F_{2,2}\) 型变量的线性变换。标准结果给出
\[
\frac{U^2}{U^2+Z_6^2}\sim \operatorname{Beta}(1,1)
\]
且 \(Z_2^2+Z_4^2\sim\chi^2_4\) 与其独立。于是
\[
V^2\sim \chi^2_2.
\]
故
\[
f_{V^2}(v)=\frac12\ee^{-v/2},\qquad v>0.
\]
这两个量都服从自由度 2 的卡方分布。
\end{solution}

\begin{problem}
某大型总体中三类特征的频率为
\[
p_1=\theta^3,\qquad p_2=3\theta(1-\theta),\qquad p_3=(1-\theta)^3,
\]
其中 \(\theta\in(0,1)\)。样本量为 \(n\)，对应观察频数为 \(N_j\)。
\begin{enumerate}[label=(\alph*)]
\item 构造近似水平 \((1-\alpha)\) 的极大似然置信集。
\item 推导频数代入估计量
\[
\hat\theta_2=1-(N_3/n)^{1/3}
\]
的渐近分布。
\end{enumerate}
\end{problem}

\begin{solution}
设样本中三类计数为 \((N_1,N_2,N_3)\)，其对数似然为
\[
\ell(\theta)=N_1\log\theta^3+N_2\log\bigl(3\theta(1-\theta)\bigr)+N_3\log(1-\theta)^3.
\]
整理得
\[
\ell(\theta)= (3N_1+N_2)\log\theta +(N_2+3N_3)\log(1-\theta)+N_2\log3.
\]
极大似然估计由
\[
\frac{\partial\ell}{\partial\theta}= \frac{3N_1+N_2}{\theta}-\frac{N_2+3N_3}{1-\theta}=0
\]
得到
\[
\hat\theta=\frac{3N_1+N_2}{3n}.
\]
\begin{enumerate}[label=(\alph*)]
\item Fisher 信息为
\[
I(\theta)=n\left(\frac{3}{\theta}+\frac{1}{1-\theta}\right)\frac13
\]
等价地可从 multinomial 似然二阶导得到
\[
I(\theta)=\frac{n}{\theta(1-\theta)}.
\]
于是大样本近似
\[
\frac{\hat\theta-\theta}{\sqrt{\theta(1-\theta)/n}}\approx N(0,1).
\]
因此近似 \((1-\alpha)\) 置信集可写为
\[
\left\{\theta:\ \left|\frac{\hat\theta-\theta}{\sqrt{\theta(1-\theta)/n}}\right|\le z_{1-\alpha/2}\right\}.
\]
用 \(\hat\theta\) 代入标准误即可得到实际可计算的 Wald 型置信集。

\item 由于
\[
\hat\theta_2=1-(N_3/n)^{1/3},
\]
且 \(N_3/n\to p_3=(1-\theta)^3\)，由 delta 方法，
\[
\sqrt n(\hat\theta_2-\theta)\xrightarrow d N(0,\sigma^2),
\]
其中
\[
\sigma^2=\left(\frac{\dd}{\dd p} (1-p^{1/3})\Big|_{p=(1-\theta)^3}\right)^2 p_3(1-p_3)
\]
或更直接由
\[
\hat\theta_2=g(\hat p_3),\qquad g(p)=1-p^{1/3}
\]
得到
\[
g'(p)=-\frac{1}{3}p^{-2/3}=-\frac{1}{3(1-\theta)^2}.
\]
由于 \(\sqrt n(\hat p_3-p_3)\to N(0,p_3(1-p_3))\)，所以
\[
\sqrt n(\hat\theta_2-\theta)\xrightarrow d N\!\left(0,\ \frac{p_3(1-p_3)}{9(1-\theta)^4}\right)
=N\!\left(0,\ \frac{\theta^3}{9(1-\theta)}\right).
\]
\end{enumerate}
\end{solution}

\begin{problem}
设
\[
S=\begin{pmatrix}\sigma&u^T\\0&S_c\end{pmatrix},
\qquad
T=\begin{pmatrix}\tau&v^T\\0&T_c\end{pmatrix},
\qquad
b=\begin{pmatrix}\beta\\b_c\end{pmatrix},
\]
其中 \(\sigma,\tau,\beta\) 为标量，\(S_c,T_c\) 为 \(n\times n\) 矩阵，\(b_c\) 为 \(n\)-维向量。设存在向量 \(x_c\) 使
\[
(S_cT_c-\lambda I)x_c=b_c
\]
且 \(w_c=T_cx_c\) 已知。证明
\[
x=\begin{pmatrix}\gamma\\x_c\end{pmatrix},
\qquad
\gamma=\frac{\beta-\sigma v^Tx_c-u^Tw_c}{\sigma\tau-\lambda}
\]
满足
\[
(ST-\lambda I)x=b.
\]
并由此导出解上三角矩阵乘积系统的 \(O(n^2)\) 与 \(O(pn^2)\) 算法。
\end{problem}

\begin{solution}
先算
\[
ST=
\begin{pmatrix}
\sigma\tau & \sigma v^T+u^TT_c\\
0 & S_cT_c
\end{pmatrix}.
\]
于是
\[
(ST-\lambda I)\begin{pmatrix}\gamma\\x_c\end{pmatrix}
=
\begin{pmatrix}
(\sigma\tau-\lambda)\gamma+(\sigma v^T+u^TT_c)x_c\\
(S_cT_c-\lambda I)x_c
\end{pmatrix}.
\]
由于 \(w_c=T_cx_c\)，上式上分量变为
\[
(\sigma\tau-\lambda)\gamma+\sigma v^Tx_c+u^Tw_c.
\]
选
\[
\gamma=\frac{\beta-\sigma v^Tx_c-u^Tw_c}{\sigma\tau-\lambda}
\]
即可使上分量等于 \(\beta\)，下分量已按假设等于 \(b_c\)。证明完成。

\paragraph{算法}
对 \(U_1U_2-\lambda I\)：
\begin{enumerate}[label=\arabic*.]
\item 先解 \((U_2-\mu I)z=b\) 或按递推把最后一层消去，得到中间向量。
\item 再解上三角系统。
\item 逐层回代，利用块三角结构把 \(n\times n\) 问题化为若干次上三角回代。
\end{enumerate}
由于上三角系统一次回代为 \(O(n^2)\)，两层乘积也是 \(O(n^2)\)。

对 \(U_1\cdots U_p-\lambda I\)，以相同的块递推思路逐层消去一个上三角因子，每层最多做一次 \(O(n^2)\) 回代，共 \(p\) 层，所以总复杂度 \(O(pn^2)\)。
\end{solution}

\begin{problem}
设 \(A\in\R^{m\times n}\)。证明存在正交矩阵 \(U\in\R^{m\times m}\)、\(V\in\R^{n\times n}\) 使
\[
U^TAV=\diag(\sigma_1,\ldots,\sigma_p),
\qquad p=\min\{m,n\},
\]
其中 \(\sigma_1\ge\cdots\ge\sigma_p\ge0\)。若 \(\rank(A)=r\)，证明对任意 \(k<r\),
\[
\min_{\rank(B)=k}\|A-B\|_2=\sigma_{k+1}.
\]
\end{problem}

\begin{solution}
令 \(A^TA\) 的特征值为
\[
\sigma_1^2\ge\cdots\ge\sigma_p^2\ge0,
\]
对应正交特征向量 \(v_i\)。定义
\[
u_i=\frac1{\sigma_i}Av_i,\qquad \sigma_i>0.
\]
则 \(\{u_i\}\) 正交归一。把它们补成正交基，得正交矩阵 \(U,V\) 使
\[
U^TAV=\diag(\sigma_1,\ldots,\sigma_p).
\]
这是奇异值分解。

对任意秩 \(k\) 矩阵 \(B\)，由 Eckart-Young-Mirsky 定理，
\[
\|A-B\|_2\ge \sigma_{k+1}.
\]
并且取
\[
A_k=\sum_{i=1}^k \sigma_i u_iv_i^T,
\]
有 \(\rank(A_k)=k\) 且
\[
\|A-A_k\|_2=\sigma_{k+1}.
\]
故最小值即为 \(\sigma_{k+1}\)。
\end{solution}

\subsection{2011 年个人赛：应用、计算、概率与统计}

\begin{problem}
设加权内积
\[
(f,g)=\int_a^b\rho(x)f(x)g(x)\dd x,\qquad \rho(x)>0.
\]
定义多项式序列
\[
p_0(x)=1,\qquad p_1(x)=x-a_1,
\]
\[
p_n(x)=(x-a_n)p_{n-1}(x)-b_np_{n-2}(x),\qquad n\ge2,
\]
其中
\[
a_n=\frac{(xp_{n-1},p_{n-1})}{(p_{n-1},p_{n-1})},\qquad
b_n=\frac{(xp_{n-1},p_{n-2})}{(p_{n-2},p_{n-2})}.
\]
\begin{enumerate}[label=(\alph*)]
\item 证明 \(\{p_n(x)\}\) 是一组正交单项式。
\item 设 \(\{q_n(x)\}\) 是与该内积对应的正交单项式序列。证明其唯一，并且它在所有 \(n\) 次首一多项式中使 \(\|q_n\|\) 最小。
\item 若 \(\rho(x)=1/\sqrt{1-x^2}\)、\([a,b]=[-1,1]\)，证明对应正交序列为 Chebyshev 多项式
\[
T_n(x)=\cos(n\arccos x)
\]
并满足递推
\[
T_{n+1}(x)=2xT_n(x)-T_{n-1}(x).
\]
\item 求权函数 \(\rho(x)=1/\sqrt{1-x^2}\) 下 \(f(x)=x^3\) 在 \([-1,1]\) 上的最佳二次逼近。
\end{enumerate}
\end{problem}

\begin{solution}
\begin{enumerate}[label=(\alph*)]
\item 由定义 \(p_n\) 为首一多项式。用递推式与 \(a_n,b_n\) 的选择证明
\[
(p_n,p_j)=0,\qquad j<n.
\]
确切地，对 \(j\le n-2\),
\[
(p_n,p_j)=(xp_{n-1},p_j)-a_n(p_{n-1},p_j)-b_n(p_{n-2},p_j)=0
\]
因为 \(xp_{n-1}\) 与低阶正交性递推可逐级下传；对 \(j=n-1\)，
\[
(p_n,p_{n-1})=(xp_{n-1},p_{n-1})-a_n(p_{n-1},p_{n-1})=0.
\]
故正交。

\item 设 \(q_n\) 与 \(r_n\) 都是首一正交 \(n\) 次多项式，则 \(q_n-r_n\) 次数 \(<n\) 且与所有低阶正交，因而与自身正交，只能为零，所以唯一。

对任意首一 \(n\) 次多项式 \(p=q_n+s\)（其中 \(\deg s<n\)），有
\[
\|p\|^2=\|q_n\|^2+\|s\|^2\ge\|q_n\|^2,
\]
因为 \(q_n\perp s\)。故 \(q_n\) 最小。

\item 取 \(x=\cos\theta\)，则
\[
T_n(x)=\cos(n\theta).
\]
利用三角恒等式
\[
\cos((n+1)\theta)=2\cos\theta\cos(n\theta)-\cos((n-1)\theta)
\]
得到递推式。又
\[
\int_{-1}^1\frac{T_m(x)T_n(x)}{\sqrt{1-x^2}}\dd x
=\int_0^\pi \cos(m\theta)\cos(n\theta)\dd\theta,
\]
所以 \(\{T_n\}\) 两两正交。

\item 在该权函数下，奇函数 \(x^3\) 的最佳二次逼近只能由奇次项 \(x\) 组成。故设
\[
P_2(x)=ax.
\]
正交条件要求
\[
\int_{-1}^1\frac{(x^3-ax)x}{\sqrt{1-x^2}}\dd x=0.
\]
利用
\[
\int_{-1}^1\frac{x^2}{\sqrt{1-x^2}}\dd x=\frac\pi2,\qquad
\int_{-1}^1\frac{x^4}{\sqrt{1-x^2}}\dd x=\frac{3\pi}{8},
\]
得
\[
 \frac{3\pi}{8}-a\frac{\pi}{2}=0,\qquad a=\frac34.
\]
因此最佳二次逼近为
\[
P_2(x)=\frac34x.
\]
\end{enumerate}
\end{solution}

\begin{problem}
两个五次多项式 \(p(x),q(x)\) 满足
\[
p(i)=q(i)=\frac1i,\qquad i=2,3,4,5,6,
\]
且
\[
p(1)=1,\qquad q(1)=2.
\]
求 \(p(0)-q(0)\)。
\end{problem}

\begin{solution}
设
\[
r(x)=p(x)-q(x).
\]
则 \(r\) 是五次以下多项式，并满足
\[
r(i)=0,\quad i=2,3,4,5,6,\qquad r(1)=-1.
\]
因此
\[
r(x)=c(x-2)(x-3)(x-4)(x-5)(x-6).
\]
代入 \(x=1\)：
\[
-1=c(-1)(-2)(-3)(-4)(-5)=-120c,
\]
故
\[
c=\frac1{120}.
\]
于是
\[
p(0)-q(0)=r(0)=\frac1{120}(-2)(-3)(-4)(-5)(-6)=6.
\]
\end{solution}

\begin{problem}
把 \(m\) 个黑球和 \(n\) 个红球放在罐中。设 \(1\le r\le k\le n\)。
\begin{enumerate}[label=(\arabic*)]
\item 每次不放回抽球时，第 \(k\) 次抽到第 \(r\) 个红球的概率是多少？
\item 每次放回抽球时，前 \(k\) 次恰好抽到奇数个红球的概率是多少？
\end{enumerate}
\end{problem}

\begin{solution}
\begin{enumerate}[label=(\arabic*)]
\item 不放回抽球时，球的排列等可能。第 \(r\) 个红球在全排列中的位置服从在 \(m+n\) 个位置中均匀分布；且要在第 \(k\) 次被抽到，等价于该球的位置正好是 \(k\)。故概率
\[
\frac1{m+n}.
\]
若题意是“第 \(k\) 次抽到的恰为第 \(r\) 个红球”，则概率为
\[
\frac{\binom{k-1}{r-1}\binom{m+n-k}{n-r}}{\binom{m+n}{n}},
\]
因为前 \(k-1\) 次中恰有 \(r-1\) 个红球。

\item 每次放回时，每次抽到红球概率为 \(p=n/(m+n)\)。前 \(k\) 次红球数 \(X\sim\operatorname{Bin}(k,p)\)。所求为
\[
P(X\ \text{为奇数})
=\sum_{j\ \text{odd}}\binom{k}{j}p^j(1-p)^{k-j}.
\]
用
\[
(1-2p)^k=\sum_{j=0}^k\binom{k}{j}(1-p)^{k-j}(-p)^j
\]
得奇数项和
\[
\frac{1-(1-2p)^k}{2}.
\]
代入 \(p=n/(m+n)\)，故答案为
\[
\frac{1-\left(1-\frac{2n}{m+n}\right)^k}{2}.
\]
\end{enumerate}
\end{solution}

\begin{problem}
设 \(X,Y\) 独立同分布。证明
\[
\E|X+Y|\ge\E|X|.
\]
\end{problem}

\begin{solution}
写
\[
X=X^+-X^-,\qquad X^\pm\ge0.
\]
由对称性可设 \(\E X\ge0\)，则 \(\E X^+\ge \E X^-\)。条件于 \(X\)，有
\[
\E(|X+Y|\mid X)=\E|X+Y|.
\]
更简单地，利用三角不等式与独立性：
\[
\E|X+Y|
\ge \E\bigl||X|-|Y|\bigr|.
\]
但此式并不直接给出结论。题目提示的标准做法是比较正负部分。因 \(Y\) 与 \(X\) 同分布，
\[
\E|X+Y|
=\E\E(|X+Y|\mid X)
\]
把 \(Y\) 平移后，函数 \(g(x)=\E|x+Y|\) 是凸函数且最小点在任一中位数处。对称同分布下，\(0\) 是 \(Y\) 的中位数之一，因此 \(g(x)\ge g(0)=\E|Y|=\E|X|\)。故
\[
\E|X+Y|\ge\E|X|.
\]
\end{solution}

\begin{problem}
设 \(X_1,\ldots,X_n\) 来自 Bernoulli(\(p_1\))，\(Y_1,\ldots,Y_n\) 为独立 Bernoulli(\(p_2\)) 样本。
\begin{enumerate}[label=(\alph*)]
\item 给出 \(\theta=p_1-p_2\) 的极小充分统计量与 UMVU 估计。
\item 给出 \(v(p_1)=p_1(1-p_1)\) 的无偏估计方差的 Cramer-Rao 下界，或其 UMVU 估计的方差下界。
\item 计算检验
\[
\sqrt n\left|\frac{\hat p_1-\hat p_2}{\sqrt{2\hat p\hat q}}\right|\ge z_{1-\alpha}
\]
在 \(p_1=p,\ p_2=p+n^{-1/2}\Delta\) 时的渐近功效，其中 \(\hat p=\frac12(\hat p_1+\hat p_2)\)。
\end{enumerate}
\end{problem}

\begin{solution}
\begin{enumerate}[label=(\alph*)]
\item 设
\[
S_X=\sum_{i=1}^nX_i,\qquad S_Y=\sum_{i=1}^nY_i.
\]
则 \((S_X,S_Y)\) 为充分统计量，且由指数族结构也是极小充分统计量。

对 \(\theta=p_1-p_2\) 的无偏估计，自然取
\[
\hat\theta=\hat p_1-\hat p_2=\frac{S_X}{n}-\frac{S_Y}{n},
\]
并用 Rao-Blackwell 化可知其即为 UMVU，因为它已是 \((S_X,S_Y)\) 的函数，且在该二参数完全充分统计下无偏函数唯一。

\item 对单个 Bernoulli(\(p_1\))，Fisher 信息为
\[
I(p_1)=\frac{n}{p_1(1-p_1)}.
\]
对 \(v(p_1)=p_1(1-p_1)\)，无偏估计的方差下界由 delta 形式 C-R 不等式给出
\[
\Var(\tilde v)\ge \frac{(v'(p_1))^2}{I(p_1)}
=\frac{(1-2p_1)^2p_1(1-p_1)}{n}.
\]
若取 UMVU 估计 \(\hat v=\frac{S_X(n-S_X)}{n(n-1)}\)，其方差为相应有限样本表达式，且满足上述下界至多渐近相等。

\item 在局部备择下，
\[
\sqrt n(\hat p_1-\hat p_2)\xrightarrow d N(\Delta,2p(1-p)).
\]
同时 \(\hat p\to p\)。所以检验统计量渐近分布为
\[
\left|\frac{N(\Delta,2p(1-p))}{\sqrt{2p(1-p)}}\right|
=|N(\Delta/\sqrt{2p(1-p)},1)|.
\]
故渐近功效为
\[
P\left(|Z+\mu|\ge z_{1-\alpha}\right),\qquad
\mu=\frac{\Delta}{\sqrt{2p(1-p)}}.
\]
\end{enumerate}
\end{solution}

\begin{problem}
设 \(\theta\) 的独立无偏测量来自两台仪器，分别满足 \(N(0,\sigma_i^2)\) 误差，\(\sigma_1^2,\sigma_2^2\) 已知。记录数据为 \((a_1,y_1),\ldots,(a_n,y_n)\)，其中 \(a_m\in\{1,2\}\) 表示所用仪器，且 \(P(a_m=1)=P(a_m=2)=1/2\)。
\begin{enumerate}[label=(\alph*)]
\item 求 MLE \(\hat\theta\)。
\item 写出期望 Fisher 信息 \(I_\theta\) 与观察 Fisher 信息 \(I_x\)，并讨论 \(I_\theta/I_x\) 当 \(n\to\infty\) 的行为。
\item 证明 \(a=\sum_{m=1}^n(2-a_m)\) 是辅助统计量，且条件方差 \(\Var(\hat\theta\mid a)=1/I_x\)。比较 \(\hat\theta\sim N(\theta,1/I_\theta)\) 与 \(\hat\theta\sim N(\theta,1/I_x)\) 两个近似。
\end{enumerate}
\end{problem}

\begin{solution}
令
\[
\sigma_{a_m}^2=
\begin{cases}
\sigma_1^2,&a_m=1,\\
\sigma_2^2,&a_m=2.
\end{cases}
\]
则条件于 \(a\) 的对数似然为
\[
\ell(\theta)
=-\frac12\sum_{m=1}^n\left[
\log(2\pi\sigma_{a_m}^2)+\frac{(y_m-\theta)^2}{\sigma_{a_m}^2}
\right].
\]
\begin{enumerate}[label=(\alph*)]
\item 对 \(\theta\) 求导并置零：
\[
\frac{\partial\ell}{\partial\theta}
=\sum_{m=1}^n\frac{y_m-\theta}{\sigma_{a_m}^2}=0,
\]
故
\[
\hat\theta
=\frac{\sum_{m=1}^n y_m/\sigma_{a_m}^2}
{\sum_{m=1}^n1/\sigma_{a_m}^2}.
\]

\item 期望 Fisher 信息为
\[
I_\theta=\E\!\left[-\frac{\partial^2\ell}{\partial\theta^2}\right]
=\sum_{m=1}^n\E\left[\frac1{\sigma_{a_m}^2}\right]
=\frac n2\left(\frac1{\sigma_1^2}+\frac1{\sigma_2^2}\right).
\]
观察信息为
\[
I_x=\sum_{m=1}^n\frac1{\sigma_{a_m}^2}.
\]
令 \(a\) 为用仪器 1 的次数，则
\[
I_x=\frac{a}{\sigma_1^2}+\frac{n-a}{\sigma_2^2}.
\]
由大数定律，\(a/n\to1/2\) 近似几乎处处，因此
\[
\frac{I_\theta}{I_x}\to1.
\]

\item 因每次仪器选择独立且 \(P(a_m=1)=1/2\)，故 \(a\sim\operatorname{Bin}(n,1/2)\)，并且与 \(\theta\) 无关，所以 \(a\) 是辅助统计量。

条件于 \(a\)，\(\hat\theta\) 是独立高斯样本的加权平均，故
\[
\Var(\hat\theta\mid a)
=\left(\sum_{m=1}^n\frac1{\sigma_{a_m}^2}\right)^{-1}
=\frac1{I_x}.
\]
在推断上，\(\hat\theta\sim N(\theta,1/I_\theta)\) 使用平均信息，适合只看无条件大样本；\(\hat\theta\sim N(\theta,1/I_x)\) 则条件于辅助统计量，更贴近给定实验安排下的真实随机性，通常更精确。若要做基于实际样本配置的条件推断，应采用后者。
\end{enumerate}
\end{solution}

\section{笔试题}
\emph{本部分没有对应源文件。}

\section{团体赛题}

\end{document}
